\chapter{Funzioni, Goniometria e Numeri Complessi}

\section{Le Funzioni Reali}
\subsection{Introduzione e definizioni}
\subsubsection{Dominio, codominio e immagine}
\textbf{Definizione:} Dati due insiemi non vuoti $A$ e $B$, una funzione (o applicazione) $f: A \to B$ è una legge che associa a ogni elemento $x \in A$ uno e un solo elemento $y \in B$.
\begin{itemize}
    \item $A$ è detto \textit{dominio} (o insieme di definizione). Nel caso di funzioni reali di variabile reale, se non specificato, si intende il dominio naturale, ovvero il massimo sottoinsieme di $\mathbb{R}$ per cui la legge ha senso.
    \item $B$ è il \textit{codominio}, l'insieme di arrivo.
    \item L'\textit{immagine} di $f$, denotata con $f(A)$ o $\text{Im}(f)$, è il sottoinsieme del codominio costituito dai valori effettivamente assunti dalla funzione: $\text{Im}(f) = \{y \in B \mid \exists x \in A \text{ tale che } f(x) = y\}$. In generale, $\text{Im}(f) \subseteq B$.
\end{itemize}

\subsubsection{Funzioni iniettive, suriettive, biunivoche}
Sia $f: A \to B$.
\begin{itemize}
    \item \textbf{Iniettiva:} Elementi distinti del dominio hanno immagini distinte. $\forall x_1, x_2 \in A, x_1 \neq x_2 \implies f(x_1) \neq f(x_2)$, equivalentemente $\forall x_1, x_2 \in A, f(x_1) = f(x_2) \implies x_1 = x_2$.
    \item \textbf{Suriettiva:} Ogni elemento del codominio è immagine di almeno un elemento del dominio. $\forall y \in B, \exists x \in A \mid f(x) = y$. (Ovvero $f(A) = B$).
    \item \textbf{Biunivoca (o Biiettiva):} Una funzione contemporaneamente iniettiva e suriettiva. Realizza una corrispondenza uno-a-uno tra gli elementi di $A$ e di $B$.
\end{itemize}

\subsubsection{Funzione inversa e funzione composta}
\textbf{Funzione inversa:} Se $f: A \to B$ è biunivoca, è ben definita la funzione inversa $f^{-1}: B \to A$ che associa a ogni $y \in B$ l'unico $x \in A$ tale che $f(x) = y$. Il grafico di $f^{-1}$ è simmetrico a quello di $f$ rispetto alla bisettrice del primo e terzo quadrante ($y=x$).
\textbf{Funzione composta:} Date $f: A \to B$ e $g: B \to C$, la funzione composta $g \circ f: A \to C$ è definita come $(g \circ f)(x) = g(f(x))$. La composizione non è commutativa ($g \circ f \neq f \circ g$).

\subsection{Proprietà notevoli}
\subsubsection{Funzioni pari e dispari (simmetrie)}
Sia $f: D \to \mathbb{R}$ con un dominio $D$ simmetrico rispetto all'origine (se $x \in D$, allora $-x \in D$).
\begin{itemize}
    \item $f$ è \textbf{pari} se $\forall x \in D, f(-x) = f(x)$. Il grafico è simmetrico rispetto all'asse $y$. (Esempio: $y=x^2$, $y=\cos x$).
    \item $f$ è \textbf{dispari} se $\forall x \in D, f(-x) = -f(x)$. Il grafico è simmetrico rispetto all'origine. (Esempio: $y=x^3$, $y=\sin x$).
\end{itemize}

\subsubsection{Funzioni crescenti, decrescenti e monotonia}
Sia $I \subseteq D$ un intervallo.
\begin{itemize}
    \item $f$ è \textbf{strettamente crescente} in $I$ se $\forall x_1, x_2 \in I$ con $x_1 < x_2$, risulta $f(x_1) < f(x_2)$.
    \item $f$ è \textbf{debolmente crescente} in $I$ se $\forall x_1, x_2 \in I$ con $x_1 < x_2$, risulta $f(x_1) \le f(x_2)$.
\end{itemize}
Definizioni analoghe, a disuguaglianza invertita per le immagini, valgono per le funzioni decrescenti. Una funzione si dice monotòna se preserva il verso del suo andamento (crescente o decrescente) su tutto il dominio. 

\subsubsection{Funzioni periodiche}
\textbf{Definizione:} Una funzione $f$ si dice periodica di periodo $T > 0$ se $\forall x \in D$, risulta $x+T \in D$ e $f(x+T) = f(x)$. Il più piccolo $T$ per cui ciò si verifica è detto \textit{periodo fondamentale}.

\section{Funzioni Esponenziali e Logaritmiche}
\subsection{La funzione esponenziale}
\subsubsection{Definizione e grafico di $y = a^x$}
\textbf{Definizione:} Fissata una base $a \in \mathbb{R}$ con $a > 0$ e $a \neq 1$, si definisce funzione esponenziale la funzione $f: \mathbb{R} \to \mathbb{R}^+$ data da $y = a^x$.
\begin{itemize}
    \item \textbf{Se $\mathbf{a > 1}$:} La funzione è strettamente crescente, $\lim_{x \to -\infty} a^x = 0^+$, $\lim_{x \to +\infty} a^x = +\infty$.
    \item \textbf{Se $\mathbf{0 < a < 1}$:} La funzione è strettamente decrescente, $\lim_{x \to -\infty} a^x = +\infty$, $\lim_{x \to +\infty} a^x = 0^+$.
\end{itemize}
In entrambi i casi, la funzione è biunivoca sul suo codominio $\mathbb{R}^+$ (dunque invertibile) e interseca l'asse $y$ nel punto $(0,1)$.

\subsubsection{Il numero di Nepero $e$}
Base privilegiata in analisi matematica per via delle sue proprietà differenziali, il numero di Nepero $e \approx 2.71828\dots$ è un numero irrazionale e trascendente. Può essere definito rigorosamente come il limite della successione $e = \lim_{n \to +\infty} \left(1+\frac{1}{n}\right)^n$. La funzione esponenziale per eccellenza è $f(x) = e^x$.

\subsection{La funzione logaritmica}
\subsubsection{Definizione e grafico di $y = \log_a(x)$}
\textbf{Definizione:} Essendo la funzione esponenziale $y=a^x$ (con $a>0, a \neq 1$) biunivoca da $\mathbb{R}$ a $\mathbb{R}^+$, essa ammette inversa. Si definisce logaritmo in base $a$ di $x$ l'esponente da dare alla base $a$ per ottenere l'argomento $x$. In simboli: $y = \log_a(x) \iff a^y = x$.
Il dominio è $x > 0$ e il codominio è $\mathbb{R}$. Il grafico interseca l'asse $x$ in $(1,0)$ ed è asintotico all'asse $y$.

\subsubsection{Proprietà dei logaritmi}
Fissati argomenti strettamente positivi, le proprietà derivano isomorficamente da quelle delle potenze:
\begin{itemize}
    \item $\log_a(x \cdot y) = \log_a(x) + \log_a(y)$
    \item $\log_a\left(\frac{x}{y}\right) = \log_a(x) - \log_a(y)$
    \item $\log_a(x^k) = k \cdot \log_a(x)$ per ogni $k \in \mathbb{R}$.
\end{itemize}

\subsubsection{Cambiamento di base}
Per transitare da una base $a$ a una base $c$, si utilizza l'identità:
\[ \log_a(x) = \frac{\log_c(x)}{\log_c(a)} \]

\subsection{Equazioni e disequazioni esponenziali}
\subsubsection{Risoluzione per basi uguali}
Sfruttando l'iniettività della funzione esponenziale, un'equazione del tipo $a^{f(x)} = a^{g(x)}$ si riconduce all'uguaglianza degli esponenti: $f(x) = g(x)$.

\subsubsection{Risoluzione mediante logaritmi o incognita ausiliaria}
\begin{itemize}
    \item Se le basi sono diverse (es. $a^{f(x)} = b^{g(x)}$), si applica il logaritmo a entrambi i membri, abbassando gli esponenti: $f(x)\ln(a) = g(x)\ln(b)$.
    \item Se l'equazione è trinomia (es. $a^{2x} + k a^x + c = 0$), si opera la sostituzione $t = a^x$, imponendo $t > 0$, per ricondursi a un'equazione algebrica di secondo grado.
\end{itemize}

\subsection{Equazioni e disequazioni logaritmiche}
\subsubsection{Condizioni di esistenza}
A causa del dominio della funzione logaritmica, prima di applicare alcun principio risolutivo è imperativo imporre le C.E.: l'argomento di \textit{ogni} logaritmo presente deve essere posto strettamente maggiore di zero.

\subsubsection{Metodi risolutivi}
\begin{itemize}
    \item \textbf{Equazioni:} Ricondotti alla forma $\log_a(f(x)) = \log_a(g(x))$, per iniettività si passa a $f(x) = g(x)$, intersecando le soluzioni con le C.E. Se si ha $\log_a(f(x)) = k$, per definizione equivale a $f(x) = a^k$.
    \item \textbf{Disequazioni:} Dalla forma $\log_a(f(x)) \gtrless \log_a(g(x))$, si passa agli argomenti prestando massima attenzione alla base.
        \begin{itemize}
            \item Se $a > 1$ (base crescente): Il verso della disuguaglianza tra gli argomenti si \textit{mantiene} ($f(x) \gtrless g(x)$).
            \item Se $0 < a < 1$ (base decrescente): Il verso della disuguaglianza tra gli argomenti si \textit{inverte}.
        \end{itemize}
\end{itemize}

\section{Goniometria}
\subsection{Funzioni goniometriche fondamentali}
\subsubsection{Circonferenza goniometrica e misura in radianti}
Il contesto naturale della goniometria è la circonferenza goniometrica, centrata nell'origine degli assi cartesiani e di raggio unitario ($x^2+y^2=1$).
\textbf{Misura in radianti:} Un angolo al centro di $1$ radiante sottende un arco di circonferenza di lunghezza pari al raggio. In goniometria, la lunghezza dell'arco sulla circonferenza unitaria fornisce la misura dell'angolo in radianti (dunque un angolo giro vale $2\pi$).

\subsubsection{Seno, coseno e identità fondamentale}
Dato un angolo orientato $\alpha$ con vertice nell'origine e un lato sul semiasse positivo delle $x$, il lato termine interseca la circonferenza goniometrica in un punto $P$. 
\begin{itemize}
    \item Il \textbf{coseno} di $\alpha$ ($\cos \alpha$) è l'ascissa del punto $P$.
    \item Il \textbf{seno} di $\alpha$ ($\sin \alpha$) è l'ordinata del punto $P$.
\end{itemize}
Dal teorema di Pitagora discende la prima identità fondamentale: $\sin^2 \alpha + \cos^2 \alpha = 1, \forall \alpha \in \mathbb{R}$.

\subsubsection{Tangente, cotangente, secante e cosecante}
\begin{itemize}
    \item $\tan \alpha = \frac{\sin \alpha}{\cos \alpha}$ (definita per $\alpha \neq \frac{\pi}{2} + k\pi$). Geometricamente, è l'ordinata del punto di intersezione tra il prolungamento del raggio vettore e la tangente geometrica in $(1,0)$.
    \item $\cot \alpha = \frac{\cos \alpha}{\sin \alpha}$ (definita per $\alpha \neq k\pi$).
    \item $\sec \alpha = \frac{1}{\cos \alpha}$, $\csc \alpha = \frac{1}{\sin \alpha}$.
\end{itemize}

\subsubsection{Grafici delle funzioni goniometriche}
Seno e coseno sono funzioni periodiche di periodo $2\pi$, a valori limitati in $[-1, 1]$. La funzione tangente è periodica di periodo $\pi$, suriettiva su $\mathbb{R}$, e presenta asintoti verticali nei punti di non definizione.

\subsection{Archi associati e formule}
\subsubsection{Angoli notevoli e archi associati}
Sfruttando le simmetrie della circonferenza (rispetto agli assi, all'origine, e alle bisettrici), è possibile ricondurre il calcolo delle funzioni goniometriche di angoli arbitrari a quelli del primo quadrante. 
\textbf{Esempio (angoli supplementari):} $\sin(\pi - \alpha) = \sin \alpha$, $\cos(\pi - \alpha) = -\cos \alpha$.

\subsubsection{Formule di addizione e sottrazione}
Fondamentali per scomporre l'argomento:
\begin{itemize}
    \item $\cos(\alpha \pm \beta) = \cos\alpha\cos\beta \mp \sin\alpha\sin\beta$
    \item $\sin(\alpha \pm \beta) = \sin\alpha\cos\beta \pm \cos\alpha\sin\beta$
\end{itemize}

\subsubsection{Formule di duplicazione e bisezione}
Sostituendo $\alpha = \beta$ nelle precedenti:
\begin{itemize}
    \item $\sin(2\alpha) = 2\sin\alpha\cos\alpha$
    \item $\cos(2\alpha) = \cos^2\alpha - \sin^2\alpha = 1 - 2\sin^2\alpha = 2\cos^2\alpha - 1$
\end{itemize}
Le formule di bisezione si ricavano ponendo $2\alpha = x$, esplicitando $\sin(x/2)$ e $\cos(x/2)$ e introducendo la radice quadrata (il cui segno dipende dal quadrante).

\subsubsection{Formule parametriche, di prostaferesi e di Werner}
\begin{itemize}
    \item \textbf{Parametriche:} Esprimono seno e coseno razionalmente in funzione del parametro $t = \tan(\alpha/2)$. $\sin\alpha = \frac{2t}{1+t^2}$, $\cos\alpha = \frac{1-t^2}{1+t^2}$.
    \item \textbf{Prostaferesi/Werner:} Trasformano somme di funzioni goniometriche in prodotti e viceversa (utili per l'integrazione e lo studio del segno). Es: $\sin p + \sin q = 2\sin\left(\frac{p+q}{2}\right)\cos\left(\frac{p-q}{2}\right)$.
\end{itemize}

\subsection{Equazioni e disequazioni goniometriche}
\subsubsection{Equazioni elementari}
Della forma $\sin x = a$, $\cos x = b$, o $\tan x = c$. 
Se $|a| \le 1$, l'equazione $\sin x = a$ ammette le soluzioni $x = \alpha + 2k\pi \lor x = (\pi - \alpha) + 2k\pi$, con $k \in \mathbb{Z}$, dove $\alpha$ è l'angolo principale (es. $\arcsin(a)$).

\subsubsection{Equazioni lineari in seno e coseno}
Della forma $a\sin x + b\cos x + c = 0$. I metodi risolutivi includono:
1. \textit{Metodo algebrico (formule parametriche)}: Sostituzione in $t=\tan(x/2)$, riducendo a un'equazione di secondo grado fratta (verificando a parte $x = \pi + 2k\pi$).
2. \textit{Metodo grafico}: Si pone $X=\cos x$ e $Y=\sin x$, intersecando la retta $aY + bX + c = 0$ con la circonferenza $X^2+Y^2=1$.
3. \textit{Metodo dell'angolo aggiunto}: Trasformazione di $a\sin x + b\cos x$ in $R\sin(x+\phi)$.

\subsubsection{Equazioni omogenee di secondo grado}
Della forma $a\sin^2 x + b\sin x\cos x + c\cos^2 x = 0$. Assumendo $\cos x \neq 0$, si dividono tutti i membri per $\cos^2 x$, ricondicendosi a un'equazione di secondo grado in $\tan x$. Se un termine noto $d \neq 0$ è presente, lo si moltiplica per l'identità fondamentale $(\sin^2 x + \cos^2 x)$ per rendere l'equazione omogenea.

\subsubsection{Disequazioni goniometriche e sistemi}
Si risolvono prima le equazioni associate. Successivamente, le soluzioni della disequazione si individuano evidenziando graficamente sulla circonferenza goniometrica gli archi appropriati o sfruttando i grafici delle funzioni sinusoidali.

\subsection{Funzioni goniometriche inverse}
\subsubsection{Arcoseno, arcocoseno, arcotangente}
Le funzioni goniometriche non sono globalmente iniettive e pertanto non invertibili su $\mathbb{R}$. È necessario restringere il loro dominio agli intervalli di biunivocità (monotonia stretta):
\begin{itemize}
    \item $y = \arcsin(x)$: inversa di $\sin x$ su $[-\pi/2, \pi/2]$. Dominio $[-1,1]$, codominio $[-\pi/2, \pi/2]$.
    \item $y = \arccos(x)$: inversa di $\cos x$ su $[0, \pi]$. Dominio $[-1,1]$, codominio $[0, \pi]$.
    \item $y = \arctan(x)$: inversa di $\tan x$ su $(-\pi/2, \pi/2)$. Dominio $\mathbb{R}$, codominio $(-\pi/2, \pi/2)$.
\end{itemize}

\section{Numeri Complessi}
\subsection{L'insieme $\mathbb{C}$}
\subsubsection{Unità immaginaria e definizione di numero complesso}
L'esigenza di risolvere equazioni algebriche non ammettenti radici reali (es. $x^2+1=0$) conduce all'estensione del campo $\mathbb{R}$. Si definisce \textit{unità immaginaria} l'elemento $i$ tale che $i^2 = -1$.
L'insieme dei numeri complessi $\mathbb{C}$ è costituito da elementi della forma $z = a + ib$, con $a, b \in \mathbb{R}$. Il campo $\mathbb{C}$ non possiede un ordinamento totale compatibile con la struttura algebrica.

\subsubsection{Forma algebrica e operazioni fondamentali}
Nella forma $z = a + ib$, $a = \text{Re}(z)$ è la parte reale e $b = \text{Im}(z)$ è la parte immaginaria.
\begin{itemize}
    \item \textbf{Addizione:} $(a+ib) + (c+id) = (a+c) + i(b+d)$.
    \item \textbf{Moltiplicazione:} $(a+ib) \cdot (c+id) = (ac-bd) + i(ad+bc)$.
\end{itemize}

\subsubsection{Complesso coniugato e modulo}
Sia $z = a + ib$.
\begin{itemize}
    \item Il \textit{complesso coniugato} è $\bar{z} = a - ib$.
    \item Il \textit{modulo} è il numero reale non negativo $|z| = \sqrt{a^2 + b^2}$.
\end{itemize}
Vale l'importante relazione $z \bar{z} = |z|^2 = a^2+b^2$. Da questa deriva il calcolo del reciproco: $z^{-1} = \frac{\bar{z}}{|z|^2}$ (per $z \neq 0$). Di conseguenza, la divisione è $\frac{z_1}{z_2} = \frac{z_1 \bar{z}_2}{|z_2|^2}$.

\subsection{Rappresentazione geometrica}
\subsubsection{Piano di Gauss}
L'insieme $\mathbb{C}$ è in biiezione col piano euclideo $\mathbb{R}^2$. Il numero $z = a+ib$ è rappresentato dal punto (affissa) di coordinate cartesiane $(a,b)$. L'asse delle ascisse è l'asse reale, l'asse delle ordinate è l'asse immaginario.

\subsubsection{Forma trigonometrica e coordinate polari}
Passando alle coordinate polari, il punto $(a,b)$ è identificato dalla distanza dall'origine $\rho = |z| \ge 0$ (modulo) e dall'angolo $\theta$ (argomento) formato col semiasse reale positivo, definito a meno di multipli di $2\pi$.
Si ha $a = \rho \cos \theta$ e $b = \rho \sin \theta$. La forma trigonometrica è:
\[ z = \rho(\cos \theta + i \sin \theta) \]

\subsection{Operazioni in forma trigonometrica ed esponenziale}
\subsubsection{Prodotto e quoziente}
Siano $z_1 = \rho_1(\cos\theta_1 + i\sin\theta_1)$ e $z_2 = \rho_2(\cos\theta_2 + i\sin\theta_2)$. Sfruttando le formule di addizione e sottrazione della goniometria si ottiene:
\begin{itemize}
    \item Prodotto: $z_1 z_2 = \rho_1 \rho_2 [\cos(\theta_1+\theta_2) + i\sin(\theta_1+\theta_2)]$. (Si moltiplicano i moduli e si sommano gli argomenti).
    \item Quoziente: $\frac{z_1}{z_2} = \frac{\rho_1}{\rho_2} [\cos(\theta_1-\theta_2) + i\sin(\theta_1-\theta_2)]$.
\end{itemize}

\subsubsection{Potenza e Formula di De Moivre}
Per induzione, il prodotto di un complesso per sé stesso porta alla formula di De Moivre:
\[ z^n = \rho^n [\cos(n\theta) + i\sin(n\theta)], \quad \forall n \in \mathbb{N} \]

\subsubsection{Radici $n$-esime di un numero complesso}
\textbf{Teorema:} Ogni numero complesso $z \neq 0$ ammette esattamente $n$ distinte radici $n$-esime in $\mathbb{C}$. Data la periodicità delle funzioni goniometriche, esse sono date dalla formula:
\[ \omega_k = \sqrt[n]{\rho} \left[ \cos\left(\frac{\theta + 2k\pi}{n}\right) + i\sin\left(\frac{\theta + 2k\pi}{n}\right) \right], \quad k = 0, 1, \dots, n-1 \]
Nel piano di Gauss, le $n$ radici individuano i vertici di un poligono regolare a $n$ lati inscritto in una circonferenza di raggio $\sqrt[n]{\rho}$.

\subsubsection{La forma esponenziale (Formula di Eulero)}
Dalla potente formula di Eulero $e^{i\theta} = \cos \theta + i \sin \theta$, si giunge alla rappresentazione compatta ed elegante del numero complesso, che semplifica radicalmente la notazione algebrica per moltiplicazioni e potenze:
\[ z = \rho e^{i\theta} \]
\textbf{Esempio (Identità di Eulero):} Valutata per $\theta = \pi$, si ottiene l'equazione $e^{i\pi} + 1 = 0$, legame profondo tra cinque delle costanti fondamentali della matematica.

\subsection{Teorema fondamentale dell'algebra (cenni)}
Il passaggio ai numeri complessi completa la struttura algebrica dei polinomi.
\textbf{Teorema fondamentale dell'algebra (D'Alembert-Gauss):} Ogni polinomio $P(z)$ a coefficienti complessi di grado $n \ge 1$ ammette almeno una radice in $\mathbb{C}$. 
\textbf{Corollario:} Per il Teorema di Ruffini, iterando la scomposizione, ogni polinomio di grado $n$ ammette esattamente $n$ radici nel campo complesso, contate con la loro molteplicità. Si afferma per questo che $\mathbb{C}$ è un campo \textit{algebricamente chiuso}.

\section{Esercizi}

\subsection{Le Funzioni Reali: Dominio, Simmetrie, Composizione e Inversa}
\begin{multicols}{2}
\begin{enumerate}
    % Dominio (10 varianti)
    \item \tagprob{Dominio}{1} Determina il dominio: $f(x) = \frac{3x - 1}{x^2 - 4}$
    \item \tagprob{Dominio}{1} Determina il dominio: $f(x) = \sqrt{x^2 - 5x + 6}$
    \item \tagprob{Dominio}{1} Determina il dominio: $f(x) = \frac{x}{\sqrt{x - 2}}$
    \item \tagprob{Dominio}{2} Determina il dominio: $f(x) = \sqrt{\frac{x + 1}{x - 3}}$
    \item \tagprob{Dominio}{2} Determina il dominio: $f(x) = \frac{\sqrt{x+1}}{\sqrt{x-3}}$ \textit{(Confronta col precedente)}
    \item \tagprob{Dominio}{2} Determina il dominio: $f(x) = \frac{1}{|x| - 2}$
    \item \tagprob{Dominio}{2} Determina il dominio: $f(x) = \sqrt{|x-1| - 3}$
    \item \tagprob{Dominio}{3} Determina il dominio: $f(x) = \sqrt{x - \sqrt{x^2 - 1}}$
    \item \tagprob{Dominio}{3} Determina il dominio: $f(x) = \frac{x^2+1}{\sqrt{4 - |x^2-5|}}$
    \item \tagprob{Dominio}{3} Determina il dominio: $f(x) = \sqrt{1 - \sqrt{2 - x}}$
    
    % Simmetrie (Pari/Dispari) (10 varianti)
    \item \tagprob{Simmetrie}{1} Stabilisci se $f(x) = x^4 - 2x^2 + 3$ è pari, dispari o né pari né dispari.
    \item \tagprob{Simmetrie}{1} Stabilisci se $f(x) = x^3 - 5x$ è pari o dispari.
    \item \tagprob{Simmetrie}{1} Stabilisci la simmetria di $f(x) = \frac{x}{x^2 + 1}$.
    \item \tagprob{Simmetrie}{2} Stabilisci la simmetria di $f(x) = |x| \cdot x^3$.
    \item \tagprob{Simmetrie}{2} Stabilisci la simmetria di $f(x) = \frac{|x-1| + |x+1|}{2}$.
    \item \tagprob{Simmetrie}{2} Stabilisci la simmetria di $f(x) = (x+1)^2$.
    \item \tagprob{Simmetrie}{2} Dimostra che $f(x) = f_1(x) \cdot f_2(x)$ è dispari se $f_1$ è pari e $f_2$ è dispari.
    \item \tagprob{Simmetrie}{3} Stabilisci la simmetria di $f(x) = \ln\left(\frac{1+x}{1-x}\right)$.
    \item \tagprob{Simmetrie}{3} Dimostra che ogni funzione definita su $\mathbb{R}$ può essere scritta come somma di una funzione pari e una dispari.
    \item \tagprob{Simmetrie}{3} Stabilisci la simmetria di $f(x) = e^x - e^{-x}$.
    
    % Funzione Inversa e Composta (10 varianti)
    \item \tagprob{Funz Composta}{1} Date $f(x)=2x+1$ e $g(x)=x^2$, calcola $(g \circ f)(x)$.
    \item \tagprob{Funz Composta}{1} Date le funzioni precedenti, calcola $(f \circ g)(x)$ e verifica che la composizione non è commutativa.
    \item \tagprob{Funz Composta}{2} Date $f(x)=\frac{1}{x-1}$ e $g(x)=\sqrt{x}$, calcola $f(g(x))$ e il suo dominio.
    \item \tagprob{Funz Composta}{2} Esprimi $h(x) = \sqrt{x^2+1}$ come composizione di due funzioni elementari.
    \item \tagprob{Funz Inversa}{1} Dimostra che $f(x) = 3x - 5$ è biunivoca e calcola $f^{-1}(x)$.
    \item \tagprob{Funz Inversa}{1} Calcola l'inversa di $f(x) = \frac{x+2}{x-1}$.
    \item \tagprob{Funz Inversa}{2} Determina l'inversa di $f(x) = x^3 + 1$.
    \item \tagprob{Funz Inversa}{2} Data $f(x) = x^2 - 4x + 3$, restringi il dominio a $[2, +\infty)$ per renderla invertibile e calcola l'inversa.
    \item \tagprob{Funz Inversa}{3} Determina l'inversa di $f(x) = \frac{e^x - e^{-x}}{2}$.
    \item \tagprob{Funz Inversa}{3} Dimostra che se $f$ è strettamente decrescente su $\mathbb{R}$, allora è invertibile.
\end{enumerate}
\end{multicols}

\subsection{Funzioni Esponenziali e Logaritmiche}
\begin{multicols}{2}
\begin{enumerate}
    % C.E. Logaritmi (10 varianti)
    \item \tagprob{C.E. Logaritmo}{1} Determina le C.E.: $f(x) = \log_2(3x - 6)$
    \item \tagprob{C.E. Logaritmo}{1} Determina le C.E.: $f(x) = \ln(4 - x^2)$
    \item \tagprob{C.E. Logaritmo}{1} Determina le C.E.: $f(x) = \log_{1/2}\left(\frac{x}{x+1}\right)$
    \item \tagprob{C.E. Logaritmo}{2} Determina le C.E.: $f(x) = \ln(x^2 + 1)$
    \item \tagprob{C.E. Logaritmo}{2} Determina le C.E.: $f(x) = \log_3(|x| - 2)$
    \item \tagprob{C.E. Logaritmo}{2} Determina le C.E.: $f(x) = \sqrt{\ln(x - 1)}$
    \item \tagprob{C.E. Logaritmo}{2} Determina le C.E.: $f(x) = \ln(\ln x)$
    \item \tagprob{C.E. Logaritmo}{3} Determina le C.E.: $f(x) = \ln(e^x - 1)$
    \item \tagprob{C.E. Logaritmo}{3} Determina le C.E.: $f(x) = \frac{1}{\log_2(x^2 - 3x + 3)}$
    \item \tagprob{C.E. Logaritmo}{3} Determina le C.E.: $f(x) = \log_x(2 - x)$ \textit{(Attenzione alla base)}
    
    % Equazioni Esponenziali e Logaritmiche (10 varianti)
    \item \tagprob{Eq Esponenziale}{1} Risolvi: $2^{3x-1} = 16$
    \item \tagprob{Eq Esponenziale}{1} Risolvi: $3^{x^2 - 4} = 1$
    \item \tagprob{Eq Esponenziale}{2} Risolvi: $5^{2x} - 6 \cdot 5^x + 5 = 0$
    \item \tagprob{Eq Esponenziale}{2} Risolvi: $e^{2x} + 2e^x - 3 = 0$
    \item \tagprob{Eq Esponenziale}{3} Risolvi tramite logaritmi: $2^{x+1} = 3^{x-1}$
    \item \tagprob{Eq Logaritmica}{1} Risolvi: $\log_2(x-1) = 3$
    \item \tagprob{Eq Logaritmica}{1} Risolvi: $\ln(2x) + \ln(3) = \ln(12)$
    \item \tagprob{Eq Logaritmica}{2} Risolvi: $\log_3(x) + \log_3(x-8) = 2$
    \item \tagprob{Eq Logaritmica}{2} Risolvi: $\log_{1/2}^2(x) - \log_{1/2}(x) - 2 = 0$
    \item \tagprob{Eq Logaritmica}{3} Risolvi cambiando base: $\log_2(x) + \log_4(x) = 3$
    
    % Disequazioni Esponenziali e Logaritmiche (10 varianti)
    \item \tagprob{Dis Esponenziale}{1} Risolvi: $3^{2x-1} \ge 27$
    \item \tagprob{Dis Esponenziale}{1} Risolvi: $\left(\frac{1}{2}\right)^{x+3} < \frac{1}{8}$
    \item \tagprob{Dis Esponenziale}{2} Risolvi: $4^x - 2^x - 2 > 0$
    \item \tagprob{Dis Esponenziale}{2} Risolvi: $e^{x^2 - 3x} < \frac{1}{e^2}$
    \item \tagprob{Dis Esponenziale}{3} Risolvi: $2^x \ge 3$
    \item \tagprob{Dis Logaritmica}{1} Risolvi: $\ln(x - 2) \ge 0$
    \item \tagprob{Dis Logaritmica}{1} Risolvi (attenzione alla base): $\log_{1/3}(2x-5) > -1$
    \item \tagprob{Dis Logaritmica}{2} Risolvi: $\log_2(x^2 - 2x) < 3$
    \item \tagprob{Dis Logaritmica}{2} Risolvi: $\ln^2(x) - 4 \ge 0$
    \item \tagprob{Dis Logaritmica}{3} Risolvi: $\log_{1/2}(x) > \log_{1/2}(x^2 - 2)$
\end{enumerate}
\end{multicols}

\subsection{Goniometria}
\begin{multicols}{2}
\begin{enumerate}
    % Identità, Archi Associati e Formule (10 varianti)
    \item \tagprob{Archi Associati}{1} Semplifica: $\sin(\pi - \alpha) + \cos(\pi/2 + \alpha) - \sin(2\pi - \alpha)$
    \item \tagprob{Archi Associati}{1} Calcola senza calcolatrice: $\cos(120^\circ) + \sin(210^\circ) + \tan(315^\circ)$
    \item \tagprob{Identità}{1} Verifica: $\tan \alpha + \cot \alpha = \frac{1}{\sin \alpha \cos \alpha}$
    \item \tagprob{Identità}{2} Verifica: $\frac{\sin^3 \alpha - \cos^3 \alpha}{\sin \alpha - \cos \alpha} = 1 + \sin \alpha \cos \alpha$
    \item \tagprob{Addizione}{1} Calcola il valore esatto di $\sin(75^\circ)$ sfruttando $75^\circ = 45^\circ + 30^\circ$.
    \item \tagprob{Addizione}{2} Sviluppa e semplifica: $\cos(\alpha + \pi/6) + \cos(\alpha - \pi/6)$
    \item \tagprob{Duplicazione}{1} Sapendo che $\sin \alpha = \frac{3}{5}$ e $0 < \alpha < \pi/2$, calcola $\sin(2\alpha)$ e $\cos(2\alpha)$.
    \item \tagprob{Bisezione}{2} Calcola il valore esatto di $\cos(\pi/8)$.
    \item \tagprob{Prostaferesi}{2} Trasforma in prodotto: $\sin(5\alpha) - \sin(3\alpha)$
    \item \tagprob{Parametriche}{3} Esprimi l'espressione $y = \frac{1 - \cos x}{\sin x}$ in funzione di $t = \tan(x/2)$ e semplifica.
    
    % Equazioni Goniometriche (10 varianti)
    \item \tagprob{Eq Goniometrica}{1} Elementare: $2\sin x - 1 = 0$
    \item \tagprob{Eq Goniometrica}{1} Elementare: $\sqrt{3}\tan x - 1 = 0$
    \item \tagprob{Eq Goniometrica}{2} Riconducibile: $2\cos^2 x - \cos x - 1 = 0$
    \item \tagprob{Eq Goniometrica}{2} Scomposizione: $\sin(2x) - \sin x = 0$
    \item \tagprob{Eq Goniometrica}{2} Lineare: $\sin x + \cos x = 1$ \textit{(Usa metodo grafico o angolo aggiunto)}
    \item \tagprob{Eq Goniometrica}{2} Lineare: $\sqrt{3}\sin x - \cos x = 2$
    \item \tagprob{Eq Goniometrica}{2} Omogenea: $\sin^2 x - 3\sin x \cos x + 2\cos^2 x = 0$
    \item \tagprob{Eq Goniometrica}{3} Omogenea non nulla: $3\sin^2 x + \sin x \cos x - 2\cos^2 x = 1$
    \item \tagprob{Eq Goniometrica}{3} Con formule: $\cos(2x) + 3\sin x - 2 = 0$
    \item \tagprob{Eq Goniometrica}{3} Con argomenti multipli: $\cos(3x) = \cos x$
    
    % Disequazioni e Funz Inverse (10 varianti)
    \item \tagprob{Dis Goniometrica}{1} Risolvi: $2\cos x - 1 \ge 0$ in $[0, 2\pi]$
    \item \tagprob{Dis Goniometrica}{1} Risolvi: $\tan x \le 1$ in $(-\pi/2, \pi/2)$
    \item \tagprob{Dis Goniometrica}{2} Risolvi: $2\sin^2 x + \sin x - 1 > 0$
    \item \tagprob{Dis Goniometrica}{2} Risolvi: $\sqrt{3}\sin x - \cos x < 0$
    \item \tagprob{Sist Goniometrico}{2} Risolvi in $[0, 2\pi]$: $\begin{cases} \sin x > 0 \\ \cos x < 1/2 \end{cases}$
    \item \tagprob{Funz Inversa}{1} Calcola: $\arcsin(1/2) + \arccos(-\sqrt{2}/2)$
    \item \tagprob{Funz Inversa}{2} Determina il dominio di $f(x) = \arcsin(2x - 3)$.
    \item \tagprob{Funz Inversa}{2} Calcola l'espressione: $\cos(\arcsin(4/5))$.
    \item \tagprob{Funz Inversa}{3} Risolvi l'equazione: $\arctan x + \arctan(2x) = \pi/4$.
    \item \tagprob{Funz Inversa}{3} Dimostra l'identità: $\arcsin x + \arccos x = \pi/2$ per ogni $x \in [-1, 1]$.
\end{enumerate}
\end{multicols}

\subsection{Numeri Complessi}
\begin{multicols}{2}
\begin{enumerate}
    % Forma algebrica e potenze di i (10 varianti)
    \item \tagprob{Algebrica}{1} Calcola $(2 - 3i) + (4 + 5i) - (1 - i)$.
    \item \tagprob{Algebrica}{1} Calcola il prodotto: $(3 + 2i)(3 - 2i)$.
    \item \tagprob{Algebrica}{1} Esegui la divisione: $\frac{1 + 2i}{2 - i}$.
    \item \tagprob{Algebrica}{2} Calcola il modulo di $z = \frac{3+4i}{1-i}$.
    \item \tagprob{Potenze di i}{1} Semplifica: $i^{15} - i^{22} + i^{33}$.
    \item \tagprob{Potenze di i}{2} Calcola $(1+i)^4$ in forma algebrica.
    \item \tagprob{Eq Algebrica}{2} Trova i numeri reali $x,y$ tali che $(x+iy)(1-i) = 3 + i$.
    \item \tagprob{Eq Algebrica}{2} Risolvi l'equazione $z + 2\bar{z} = 3 - i$.
    \item \tagprob{Eq Algebrica}{3} Risolvi l'equazione $z \cdot \bar{z} - i z = 2 + 2i$.
    \item \tagprob{Luoghi Geom}{3} Disegna nel piano di Gauss l'insieme dei punti $z$ tali che $|z - 2i| \le 3$.
    
    % Forma Trigonometrica, Esponenziale e Radici (10 varianti)
    \item \tagprob{Trigonometrica}{1} Scrivi in forma trigonometrica: $z = 1 - \sqrt{3}i$.
    \item \tagprob{Trigonometrica}{1} Scrivi in forma trigonometrica: $z = -2i$.
    \item \tagprob{Trigonometrica}{2} Calcola il prodotto in forma trigonometrica tra $z_1 = 2(\cos(\pi/6) + i\sin(\pi/6))$ e $z_2 = 3(\cos(\pi/3) + i\sin(\pi/3))$.
    \item \tagprob{Esponenziale}{1} Scrivi in forma esponenziale: $z = -1 + i$.
    \item \tagprob{Esponenziale}{2} Scrivi in forma algebrica il numero $z = 2e^{i\frac{3\pi}{4}}$.
    \item \tagprob{De Moivre}{2} Usando De Moivre, calcola $(\sqrt{3} + i)^6$.
    \item \tagprob{Radici C}{2} Calcola le 3 radici cubiche dell'unità ($z^3 = 1$) e rappresentale nel piano di Gauss.
    \item \tagprob{Radici C}{3} Calcola le radici quarte di $z = -16$ ed esprimile in forma algebrica.
    \item \tagprob{Eq in C}{2} Risolvi nel campo complesso: $x^2 - 4x + 13 = 0$.
    \item \tagprob{Eq in C}{3} Risolvi nel campo complesso: $z^4 + 3z^2 - 4 = 0$ \textit{(Biquadratica)}.
\end{enumerate}
\end{multicols}